\textbf{COLUNI 2025 -- Questão 06}\par\medskip Leia o segmento abaixo.

\par

Belonísia se sentiu uma sombra de Bibiana durante aqueles dias. Havia se esquivado a vida toda daquele papel, desde que, de forma quase instintiva, a irmã passou a falar por ela. Desde que permitiu que Bibiana conhecesse seus sentimentos mais íntimos. Da mesma forma, se apossava do que se movimentava feroz no pensamento da irmã. Se sentia, mais que nunca, unida pelo que parecia ser um destino inevitável a se traçar nas trilhas de suas vidas. Passado tanto tempo, não era mais preciso nenhuma comunicação visível, seja pela troca de olhares ou pela leitura dos gestos. O ar, sentia, poderia vibrar de forma involuntária transmitindo o mal-estar físico e mental que a outra emanava. Poderia transmitir suas agitações e suas vontades. Esses dias foram cruciais para que percebesse o quanto estavam ajustadas em suas compreensões. Belonísia havia desenvolvido essa percepção expandida em relação às pessoas, mais ainda quando se referia à irmã, a sua voz no mundo onde se movimentava em silêncio. (VIEIRA JUNIOR, Itamar. \textit{Torto arado}. São Paulo: Todavia, 2019. p. 217)

\par

De acordo com o segmento, assinale a afirmativa INCORRETA.\begin{enumerate}
\item Em: “\underline{Belonísia} havia desenvolvido essa percepção expandida em relação às pessoas [...]”, o termo
       destacado acentua-se por respeitar a regra das proparoxítonas.
\item Em: “[...] a vida toda daquele papel, \underline{desde que}, de forma quase instintiva, a irmã passou a falar por ela
       [...]”, a expressão destacada estabelece uma ideia de tempo.
\item Em: “[...] unida pelo que parecia ser um destino \underline{inevitável} a se traçar [...]”, o termo destacado poderia ser
       substituído, sem prejuízo de sentido, por implacável.
\item Em: “Esses dias foram cruciais para que \underline{percebesse} o quanto estavam ajustadas [...]”, o termo destacado
       está flexionado no modo subjuntivo.\end{enumerate}
